\section{Cardiac MRI}


\section*{Patient focus}
This chapter explains the examination or clinical problem in accessible language. Patients should use it to prepare questions and understand the purpose, experience, benefits, and limitations of imaging; local instructions and the official report take priority.


\subsubsection*{What Is This Test?}
Cardiac MRI provides the most detailed images of heart structure and function without radiation. It excels at tissue characterization (scar, inflammation, infiltration) and is the gold standard for heart muscle diseases.

\subsubsection*{Alternative Names}
Cardiac MRI, cardiac magnetic resonance, CMR, cardiovascular MRI


\subsubsection*{Principle}
Cardiac MRI uses ECG-gated sequences to freeze cardiac motion. Cine imaging (balanced SSFP) visualizes ventricular function and wall motion. Late gadolinium enhancement (LGE) imaging identifies myocardial scar and fibrosis. T1 and T2 mapping quantify tissue characteristics.

\subsubsection*{History}
Cardiac magnetic resonance imaging emerged from the development of MRI by Paul Lauterbur and Peter Mansfield in the 1970s. Dedicated cardiac protocols were refined in the 1990s, offering comprehensive assessment of heart structure and function.

\subsubsection*{Why This Test Is Ordered}
Your doctor orders a cardiac MRI when echocardiography does not provide enough detail, when you have a cardiomyopathy (heart muscle disease) that needs characterization, when myocarditis (heart muscle inflammation) is suspected, or when your doctor needs to evaluate heart muscle viability before deciding whether bypass surgery would help. It answers the question: what is causing your heart muscle problem at a tissue level?

\subsubsection*{Is This a Primary or Secondary Test}
A secondary test. Typically performed after echocardiography when more detailed information about heart muscle structure is needed.

\subsubsection*{Cardiac MRI: performance depends on the tissue question}
Cardiac MRI is interpreted against a specific suspected process, such as scar, edema, or infiltrative disease. Criteria and accuracy differ across these questions, so a pooled success rate would not be informative.

\subsubsection*{How This Test Is Performed}
You lie on a table that slides into an MRI machine, similar to a regular MRI but with additional equipment to monitor your heart rhythm. The scan takes forty-five to ninety minutes. You receive contrast dye (gadolinium) through an IV to highlight areas of heart muscle damage.

\subsubsection*{What Preparation Is Needed}
Same as for regular MRI. Tell the staff about any metal implants. Check kidney function before gadolinium contrast.

\subsubsection*{Contraindications and Precautions}
Cardiac MRI requires device-specific MR screening, ECG-gating considerations, and review of any planned contrast agent. Follow the facility's current implant and contrast protocol. Gadolinium contrast requires adequate kidney function.

\subsubsection*{Risks}
Cardiac MRI uses no ionizing radiation. Risks include: discomfort from lying in the scanner for 45--90 minutes, loud noise (ear protection is provided), and (rarely) allergic reaction to gadolinium contrast dye. Patients with certain metal implants may be at risk from the magnetic field.

\subsubsection*{Follow-up Tests}
If the cardiac MRI is abnormal, your doctor may order coronary angiography to assess for blockages, a repeat cardiac MRI to monitor changes, or a biopsy for tissue diagnosis.
\clearpage
